\begin{abstract}
Activation sparsity is often reported at one tensor, yet its computational
meaning depends on where zero operands enter the entire transformer graph.  We
define \emph{model-wide zero-product opportunity}, $\Rmodel$: the fraction of
scalar products in six declared transformer-block operations whose measured
operands contain an exact zero, with a dense language-model head retained in
the denominator.  The metric includes post-RoPE query--key and causal
probability--value products and is a logical count, not a runtime claim.  We
apply it to randomly initialized Pythia 14M, 70M, and 410M models trained for
one matched 1.493B-token MiniPile pass.  We compare evaluation-only clipping
with two trained gate topologies using an Adam-aware orthogonal $\ell_1$
pressure.  The same local zero rate can imply very different $\Rmodel$ because
operations have unequal footprints and zeros can be reused across attention
products.  At high gate threshold, the seven-site recipe reaches 27.5\%,
40.6\%, and 80.6\% measured opportunity at 14M, 70M, and 410M, respectively.
The 410M cohort also reverses the small-model trained loss response: stronger
gating improves rather than degrades validation loss, even though matched
post-hoc clipping still degrades it monotonically.  A learning-rate screen
does not repair its weak baseline, while its token exposure is only 3.7
tokens per parameter.  We therefore treat 410M as a boundary result, not a
scaling claim.  These results support model-wide accounting as the comparison
layer for sparse interventions and show why local sparsity, logical
opportunity, and realized acceleration must remain separate quantities.
\end{abstract}
